\documentclass[11pt]{article}

    \usepackage[breakable]{tcolorbox}
    \usepackage{parskip} % Stop auto-indenting (to mimic markdown behaviour)
    

    % Basic figure setup, for now with no caption control since it's done
    % automatically by Pandoc (which extracts ![](path) syntax from Markdown).
    \usepackage{graphicx}
    % Keep aspect ratio if custom image width or height is specified
    \setkeys{Gin}{keepaspectratio}
    % Maintain compatibility with old templates. Remove in nbconvert 6.0
    \let\Oldincludegraphics\includegraphics
    % Ensure that by default, figures have no caption (until we provide a
    % proper Figure object with a Caption API and a way to capture that
    % in the conversion process - todo).
    \usepackage{caption}
    \DeclareCaptionFormat{nocaption}{}
    \captionsetup{format=nocaption,aboveskip=0pt,belowskip=0pt}

    \usepackage{float}
    \floatplacement{figure}{H} % forces figures to be placed at the correct location
    \usepackage{xcolor} % Allow colors to be defined
    \usepackage{enumerate} % Needed for markdown enumerations to work
    \usepackage{geometry} % Used to adjust the document margins
    \usepackage{amsmath} % Equations
    \usepackage{amssymb} % Equations
    \usepackage{textcomp} % defines textquotesingle
    % Hack from http://tex.stackexchange.com/a/47451/13684:
    \AtBeginDocument{%
        \def\PYZsq{\textquotesingle}% Upright quotes in Pygmentized code
    }
    \usepackage{upquote} % Upright quotes for verbatim code
    \usepackage{eurosym} % defines \euro

    \usepackage{iftex}
    \ifPDFTeX
        \usepackage[T1]{fontenc}
        \IfFileExists{alphabeta.sty}{
              \usepackage{alphabeta}
          }{
              \usepackage[mathletters]{ucs}
              \usepackage[utf8x]{inputenc}
          }
    \else
        \usepackage{fontspec}
        \usepackage{unicode-math}
    \fi

    \usepackage{fancyvrb} % verbatim replacement that allows latex
    \usepackage{grffile} % extends the file name processing of package graphics
                         % to support a larger range
    \makeatletter % fix for old versions of grffile with XeLaTeX
    \@ifpackagelater{grffile}{2019/11/01}
    {
      % Do nothing on new versions
    }
    {
      \def\Gread@@xetex#1{%
        \IfFileExists{"\Gin@base".bb}%
        {\Gread@eps{\Gin@base.bb}}%
        {\Gread@@xetex@aux#1}%
      }
    }
    \makeatother
    \usepackage[Export]{adjustbox} % Used to constrain images to a maximum size
    \adjustboxset{max size={0.9\linewidth}{0.9\paperheight}}

    % The hyperref package gives us a pdf with properly built
    % internal navigation ('pdf bookmarks' for the table of contents,
    % internal cross-reference links, web links for URLs, etc.)
    \usepackage{hyperref}
    % The default LaTeX title has an obnoxious amount of whitespace. By default,
    % titling removes some of it. It also provides customization options.
    \usepackage{titling}
    \usepackage{longtable} % longtable support required by pandoc >1.10
    \usepackage{booktabs}  % table support for pandoc > 1.12.2
    \usepackage{array}     % table support for pandoc >= 2.11.3
    \usepackage{calc}      % table minipage width calculation for pandoc >= 2.11.1
    \usepackage[inline]{enumitem} % IRkernel/repr support (it uses the enumerate* environment)
    \usepackage[normalem]{ulem} % ulem is needed to support strikethroughs (\sout)
                                % normalem makes italics be italics, not underlines
    \usepackage{soul}      % strikethrough (\st) support for pandoc >= 3.0.0
    \usepackage{mathrsfs}
    

    
    % Colors for the hyperref package
    \definecolor{urlcolor}{rgb}{0,.145,.698}
    \definecolor{linkcolor}{rgb}{.71,0.21,0.01}
    \definecolor{citecolor}{rgb}{.12,.54,.11}

    % ANSI colors
    \definecolor{ansi-black}{HTML}{3E424D}
    \definecolor{ansi-black-intense}{HTML}{282C36}
    \definecolor{ansi-red}{HTML}{E75C58}
    \definecolor{ansi-red-intense}{HTML}{B22B31}
    \definecolor{ansi-green}{HTML}{00A250}
    \definecolor{ansi-green-intense}{HTML}{007427}
    \definecolor{ansi-yellow}{HTML}{DDB62B}
    \definecolor{ansi-yellow-intense}{HTML}{B27D12}
    \definecolor{ansi-blue}{HTML}{208FFB}
    \definecolor{ansi-blue-intense}{HTML}{0065CA}
    \definecolor{ansi-magenta}{HTML}{D160C4}
    \definecolor{ansi-magenta-intense}{HTML}{A03196}
    \definecolor{ansi-cyan}{HTML}{60C6C8}
    \definecolor{ansi-cyan-intense}{HTML}{258F8F}
    \definecolor{ansi-white}{HTML}{C5C1B4}
    \definecolor{ansi-white-intense}{HTML}{A1A6B2}
    \definecolor{ansi-default-inverse-fg}{HTML}{FFFFFF}
    \definecolor{ansi-default-inverse-bg}{HTML}{000000}

    % common color for the border for error outputs.
    \definecolor{outerrorbackground}{HTML}{FFDFDF}

    % commands and environments needed by pandoc snippets
    % extracted from the output of `pandoc -s`
    \providecommand{\tightlist}{%
      \setlength{\itemsep}{0pt}\setlength{\parskip}{0pt}}
    \DefineVerbatimEnvironment{Highlighting}{Verbatim}{commandchars=\\\{\}}
    % Add ',fontsize=\small' for more characters per line
    \newenvironment{Shaded}{}{}
    \newcommand{\KeywordTok}[1]{\textcolor[rgb]{0.00,0.44,0.13}{\textbf{{#1}}}}
    \newcommand{\DataTypeTok}[1]{\textcolor[rgb]{0.56,0.13,0.00}{{#1}}}
    \newcommand{\DecValTok}[1]{\textcolor[rgb]{0.25,0.63,0.44}{{#1}}}
    \newcommand{\BaseNTok}[1]{\textcolor[rgb]{0.25,0.63,0.44}{{#1}}}
    \newcommand{\FloatTok}[1]{\textcolor[rgb]{0.25,0.63,0.44}{{#1}}}
    \newcommand{\CharTok}[1]{\textcolor[rgb]{0.25,0.44,0.63}{{#1}}}
    \newcommand{\StringTok}[1]{\textcolor[rgb]{0.25,0.44,0.63}{{#1}}}
    \newcommand{\CommentTok}[1]{\textcolor[rgb]{0.38,0.63,0.69}{\textit{{#1}}}}
    \newcommand{\OtherTok}[1]{\textcolor[rgb]{0.00,0.44,0.13}{{#1}}}
    \newcommand{\AlertTok}[1]{\textcolor[rgb]{1.00,0.00,0.00}{\textbf{{#1}}}}
    \newcommand{\FunctionTok}[1]{\textcolor[rgb]{0.02,0.16,0.49}{{#1}}}
    \newcommand{\RegionMarkerTok}[1]{{#1}}
    \newcommand{\ErrorTok}[1]{\textcolor[rgb]{1.00,0.00,0.00}{\textbf{{#1}}}}
    \newcommand{\NormalTok}[1]{{#1}}

    % Additional commands for more recent versions of Pandoc
    \newcommand{\ConstantTok}[1]{\textcolor[rgb]{0.53,0.00,0.00}{{#1}}}
    \newcommand{\SpecialCharTok}[1]{\textcolor[rgb]{0.25,0.44,0.63}{{#1}}}
    \newcommand{\VerbatimStringTok}[1]{\textcolor[rgb]{0.25,0.44,0.63}{{#1}}}
    \newcommand{\SpecialStringTok}[1]{\textcolor[rgb]{0.73,0.40,0.53}{{#1}}}
    \newcommand{\ImportTok}[1]{{#1}}
    \newcommand{\DocumentationTok}[1]{\textcolor[rgb]{0.73,0.13,0.13}{\textit{{#1}}}}
    \newcommand{\AnnotationTok}[1]{\textcolor[rgb]{0.38,0.63,0.69}{\textbf{\textit{{#1}}}}}
    \newcommand{\CommentVarTok}[1]{\textcolor[rgb]{0.38,0.63,0.69}{\textbf{\textit{{#1}}}}}
    \newcommand{\VariableTok}[1]{\textcolor[rgb]{0.10,0.09,0.49}{{#1}}}
    \newcommand{\ControlFlowTok}[1]{\textcolor[rgb]{0.00,0.44,0.13}{\textbf{{#1}}}}
    \newcommand{\OperatorTok}[1]{\textcolor[rgb]{0.40,0.40,0.40}{{#1}}}
    \newcommand{\BuiltInTok}[1]{{#1}}
    \newcommand{\ExtensionTok}[1]{{#1}}
    \newcommand{\PreprocessorTok}[1]{\textcolor[rgb]{0.74,0.48,0.00}{{#1}}}
    \newcommand{\AttributeTok}[1]{\textcolor[rgb]{0.49,0.56,0.16}{{#1}}}
    \newcommand{\InformationTok}[1]{\textcolor[rgb]{0.38,0.63,0.69}{\textbf{\textit{{#1}}}}}
    \newcommand{\WarningTok}[1]{\textcolor[rgb]{0.38,0.63,0.69}{\textbf{\textit{{#1}}}}}
    \makeatletter
    \newsavebox\pandoc@box
    \newcommand*\pandocbounded[1]{%
      \sbox\pandoc@box{#1}%
      % scaling factors for width and height
      \Gscale@div\@tempa\textheight{\dimexpr\ht\pandoc@box+\dp\pandoc@box\relax}%
      \Gscale@div\@tempb\linewidth{\wd\pandoc@box}%
      % select the smaller of both
      \ifdim\@tempb\p@<\@tempa\p@
        \let\@tempa\@tempb
      \fi
      % scaling accordingly (\@tempa < 1)
      \ifdim\@tempa\p@<\p@
        \scalebox{\@tempa}{\usebox\pandoc@box}%
      % scaling not needed, use as it is
      \else
        \usebox{\pandoc@box}%
      \fi
    }
    \makeatother

    % Define a nice break command that doesn't care if a line doesn't already
    % exist.
    \def\br{\hspace*{\fill} \\* }
    % Math Jax compatibility definitions
    \def\gt{>}
    \def\lt{<}
    \let\Oldtex\TeX
    \let\Oldlatex\LaTeX
    \renewcommand{\TeX}{\textrm{\Oldtex}}
    \renewcommand{\LaTeX}{\textrm{\Oldlatex}}
    % Document parameters
    % Document title
    \title{Homework\_10}
    
    
    
    
    
    
    
% Pygments definitions
\makeatletter
\def\PY@reset{\let\PY@it=\relax \let\PY@bf=\relax%
    \let\PY@ul=\relax \let\PY@tc=\relax%
    \let\PY@bc=\relax \let\PY@ff=\relax}
\def\PY@tok#1{\csname PY@tok@#1\endcsname}
\def\PY@toks#1+{\ifx\relax#1\empty\else%
    \PY@tok{#1}\expandafter\PY@toks\fi}
\def\PY@do#1{\PY@bc{\PY@tc{\PY@ul{%
    \PY@it{\PY@bf{\PY@ff{#1}}}}}}}
\def\PY#1#2{\PY@reset\PY@toks#1+\relax+\PY@do{#2}}

\@namedef{PY@tok@w}{\def\PY@tc##1{\textcolor[rgb]{0.73,0.73,0.73}{##1}}}
\@namedef{PY@tok@c}{\let\PY@it=\textit\def\PY@tc##1{\textcolor[rgb]{0.24,0.48,0.48}{##1}}}
\@namedef{PY@tok@cp}{\def\PY@tc##1{\textcolor[rgb]{0.61,0.40,0.00}{##1}}}
\@namedef{PY@tok@k}{\let\PY@bf=\textbf\def\PY@tc##1{\textcolor[rgb]{0.00,0.50,0.00}{##1}}}
\@namedef{PY@tok@kp}{\def\PY@tc##1{\textcolor[rgb]{0.00,0.50,0.00}{##1}}}
\@namedef{PY@tok@kt}{\def\PY@tc##1{\textcolor[rgb]{0.69,0.00,0.25}{##1}}}
\@namedef{PY@tok@o}{\def\PY@tc##1{\textcolor[rgb]{0.40,0.40,0.40}{##1}}}
\@namedef{PY@tok@ow}{\let\PY@bf=\textbf\def\PY@tc##1{\textcolor[rgb]{0.67,0.13,1.00}{##1}}}
\@namedef{PY@tok@nb}{\def\PY@tc##1{\textcolor[rgb]{0.00,0.50,0.00}{##1}}}
\@namedef{PY@tok@nf}{\def\PY@tc##1{\textcolor[rgb]{0.00,0.00,1.00}{##1}}}
\@namedef{PY@tok@nc}{\let\PY@bf=\textbf\def\PY@tc##1{\textcolor[rgb]{0.00,0.00,1.00}{##1}}}
\@namedef{PY@tok@nn}{\let\PY@bf=\textbf\def\PY@tc##1{\textcolor[rgb]{0.00,0.00,1.00}{##1}}}
\@namedef{PY@tok@ne}{\let\PY@bf=\textbf\def\PY@tc##1{\textcolor[rgb]{0.80,0.25,0.22}{##1}}}
\@namedef{PY@tok@nv}{\def\PY@tc##1{\textcolor[rgb]{0.10,0.09,0.49}{##1}}}
\@namedef{PY@tok@no}{\def\PY@tc##1{\textcolor[rgb]{0.53,0.00,0.00}{##1}}}
\@namedef{PY@tok@nl}{\def\PY@tc##1{\textcolor[rgb]{0.46,0.46,0.00}{##1}}}
\@namedef{PY@tok@ni}{\let\PY@bf=\textbf\def\PY@tc##1{\textcolor[rgb]{0.44,0.44,0.44}{##1}}}
\@namedef{PY@tok@na}{\def\PY@tc##1{\textcolor[rgb]{0.41,0.47,0.13}{##1}}}
\@namedef{PY@tok@nt}{\let\PY@bf=\textbf\def\PY@tc##1{\textcolor[rgb]{0.00,0.50,0.00}{##1}}}
\@namedef{PY@tok@nd}{\def\PY@tc##1{\textcolor[rgb]{0.67,0.13,1.00}{##1}}}
\@namedef{PY@tok@s}{\def\PY@tc##1{\textcolor[rgb]{0.73,0.13,0.13}{##1}}}
\@namedef{PY@tok@sd}{\let\PY@it=\textit\def\PY@tc##1{\textcolor[rgb]{0.73,0.13,0.13}{##1}}}
\@namedef{PY@tok@si}{\let\PY@bf=\textbf\def\PY@tc##1{\textcolor[rgb]{0.64,0.35,0.47}{##1}}}
\@namedef{PY@tok@se}{\let\PY@bf=\textbf\def\PY@tc##1{\textcolor[rgb]{0.67,0.36,0.12}{##1}}}
\@namedef{PY@tok@sr}{\def\PY@tc##1{\textcolor[rgb]{0.64,0.35,0.47}{##1}}}
\@namedef{PY@tok@ss}{\def\PY@tc##1{\textcolor[rgb]{0.10,0.09,0.49}{##1}}}
\@namedef{PY@tok@sx}{\def\PY@tc##1{\textcolor[rgb]{0.00,0.50,0.00}{##1}}}
\@namedef{PY@tok@m}{\def\PY@tc##1{\textcolor[rgb]{0.40,0.40,0.40}{##1}}}
\@namedef{PY@tok@gh}{\let\PY@bf=\textbf\def\PY@tc##1{\textcolor[rgb]{0.00,0.00,0.50}{##1}}}
\@namedef{PY@tok@gu}{\let\PY@bf=\textbf\def\PY@tc##1{\textcolor[rgb]{0.50,0.00,0.50}{##1}}}
\@namedef{PY@tok@gd}{\def\PY@tc##1{\textcolor[rgb]{0.63,0.00,0.00}{##1}}}
\@namedef{PY@tok@gi}{\def\PY@tc##1{\textcolor[rgb]{0.00,0.52,0.00}{##1}}}
\@namedef{PY@tok@gr}{\def\PY@tc##1{\textcolor[rgb]{0.89,0.00,0.00}{##1}}}
\@namedef{PY@tok@ge}{\let\PY@it=\textit}
\@namedef{PY@tok@gs}{\let\PY@bf=\textbf}
\@namedef{PY@tok@ges}{\let\PY@bf=\textbf\let\PY@it=\textit}
\@namedef{PY@tok@gp}{\let\PY@bf=\textbf\def\PY@tc##1{\textcolor[rgb]{0.00,0.00,0.50}{##1}}}
\@namedef{PY@tok@go}{\def\PY@tc##1{\textcolor[rgb]{0.44,0.44,0.44}{##1}}}
\@namedef{PY@tok@gt}{\def\PY@tc##1{\textcolor[rgb]{0.00,0.27,0.87}{##1}}}
\@namedef{PY@tok@err}{\def\PY@bc##1{{\setlength{\fboxsep}{\string -\fboxrule}\fcolorbox[rgb]{1.00,0.00,0.00}{1,1,1}{\strut ##1}}}}
\@namedef{PY@tok@kc}{\let\PY@bf=\textbf\def\PY@tc##1{\textcolor[rgb]{0.00,0.50,0.00}{##1}}}
\@namedef{PY@tok@kd}{\let\PY@bf=\textbf\def\PY@tc##1{\textcolor[rgb]{0.00,0.50,0.00}{##1}}}
\@namedef{PY@tok@kn}{\let\PY@bf=\textbf\def\PY@tc##1{\textcolor[rgb]{0.00,0.50,0.00}{##1}}}
\@namedef{PY@tok@kr}{\let\PY@bf=\textbf\def\PY@tc##1{\textcolor[rgb]{0.00,0.50,0.00}{##1}}}
\@namedef{PY@tok@bp}{\def\PY@tc##1{\textcolor[rgb]{0.00,0.50,0.00}{##1}}}
\@namedef{PY@tok@fm}{\def\PY@tc##1{\textcolor[rgb]{0.00,0.00,1.00}{##1}}}
\@namedef{PY@tok@vc}{\def\PY@tc##1{\textcolor[rgb]{0.10,0.09,0.49}{##1}}}
\@namedef{PY@tok@vg}{\def\PY@tc##1{\textcolor[rgb]{0.10,0.09,0.49}{##1}}}
\@namedef{PY@tok@vi}{\def\PY@tc##1{\textcolor[rgb]{0.10,0.09,0.49}{##1}}}
\@namedef{PY@tok@vm}{\def\PY@tc##1{\textcolor[rgb]{0.10,0.09,0.49}{##1}}}
\@namedef{PY@tok@sa}{\def\PY@tc##1{\textcolor[rgb]{0.73,0.13,0.13}{##1}}}
\@namedef{PY@tok@sb}{\def\PY@tc##1{\textcolor[rgb]{0.73,0.13,0.13}{##1}}}
\@namedef{PY@tok@sc}{\def\PY@tc##1{\textcolor[rgb]{0.73,0.13,0.13}{##1}}}
\@namedef{PY@tok@dl}{\def\PY@tc##1{\textcolor[rgb]{0.73,0.13,0.13}{##1}}}
\@namedef{PY@tok@s2}{\def\PY@tc##1{\textcolor[rgb]{0.73,0.13,0.13}{##1}}}
\@namedef{PY@tok@sh}{\def\PY@tc##1{\textcolor[rgb]{0.73,0.13,0.13}{##1}}}
\@namedef{PY@tok@s1}{\def\PY@tc##1{\textcolor[rgb]{0.73,0.13,0.13}{##1}}}
\@namedef{PY@tok@mb}{\def\PY@tc##1{\textcolor[rgb]{0.40,0.40,0.40}{##1}}}
\@namedef{PY@tok@mf}{\def\PY@tc##1{\textcolor[rgb]{0.40,0.40,0.40}{##1}}}
\@namedef{PY@tok@mh}{\def\PY@tc##1{\textcolor[rgb]{0.40,0.40,0.40}{##1}}}
\@namedef{PY@tok@mi}{\def\PY@tc##1{\textcolor[rgb]{0.40,0.40,0.40}{##1}}}
\@namedef{PY@tok@il}{\def\PY@tc##1{\textcolor[rgb]{0.40,0.40,0.40}{##1}}}
\@namedef{PY@tok@mo}{\def\PY@tc##1{\textcolor[rgb]{0.40,0.40,0.40}{##1}}}
\@namedef{PY@tok@ch}{\let\PY@it=\textit\def\PY@tc##1{\textcolor[rgb]{0.24,0.48,0.48}{##1}}}
\@namedef{PY@tok@cm}{\let\PY@it=\textit\def\PY@tc##1{\textcolor[rgb]{0.24,0.48,0.48}{##1}}}
\@namedef{PY@tok@cpf}{\let\PY@it=\textit\def\PY@tc##1{\textcolor[rgb]{0.24,0.48,0.48}{##1}}}
\@namedef{PY@tok@c1}{\let\PY@it=\textit\def\PY@tc##1{\textcolor[rgb]{0.24,0.48,0.48}{##1}}}
\@namedef{PY@tok@cs}{\let\PY@it=\textit\def\PY@tc##1{\textcolor[rgb]{0.24,0.48,0.48}{##1}}}

\def\PYZbs{\char`\\}
\def\PYZus{\char`\_}
\def\PYZob{\char`\{}
\def\PYZcb{\char`\}}
\def\PYZca{\char`\^}
\def\PYZam{\char`\&}
\def\PYZlt{\char`\<}
\def\PYZgt{\char`\>}
\def\PYZsh{\char`\#}
\def\PYZpc{\char`\%}
\def\PYZdl{\char`\$}
\def\PYZhy{\char`\-}
\def\PYZsq{\char`\'}
\def\PYZdq{\char`\"}
\def\PYZti{\char`\~}
% for compatibility with earlier versions
\def\PYZat{@}
\def\PYZlb{[}
\def\PYZrb{]}
\makeatother


    % For linebreaks inside Verbatim environment from package fancyvrb.
    \makeatletter
        \newbox\Wrappedcontinuationbox
        \newbox\Wrappedvisiblespacebox
        \newcommand*\Wrappedvisiblespace {\textcolor{red}{\textvisiblespace}}
        \newcommand*\Wrappedcontinuationsymbol {\textcolor{red}{\llap{\tiny$\m@th\hookrightarrow$}}}
        \newcommand*\Wrappedcontinuationindent {3ex }
        \newcommand*\Wrappedafterbreak {\kern\Wrappedcontinuationindent\copy\Wrappedcontinuationbox}
        % Take advantage of the already applied Pygments mark-up to insert
        % potential linebreaks for TeX processing.
        %        {, <, #, %, $, ' and ": go to next line.
        %        _, }, ^, &, >, - and ~: stay at end of broken line.
        % Use of \textquotesingle for straight quote.
        \newcommand*\Wrappedbreaksatspecials {%
            \def\PYGZus{\discretionary{\char`\_}{\Wrappedafterbreak}{\char`\_}}%
            \def\PYGZob{\discretionary{}{\Wrappedafterbreak\char`\{}{\char`\{}}%
            \def\PYGZcb{\discretionary{\char`\}}{\Wrappedafterbreak}{\char`\}}}%
            \def\PYGZca{\discretionary{\char`\^}{\Wrappedafterbreak}{\char`\^}}%
            \def\PYGZam{\discretionary{\char`\&}{\Wrappedafterbreak}{\char`\&}}%
            \def\PYGZlt{\discretionary{}{\Wrappedafterbreak\char`\<}{\char`\<}}%
            \def\PYGZgt{\discretionary{\char`\>}{\Wrappedafterbreak}{\char`\>}}%
            \def\PYGZsh{\discretionary{}{\Wrappedafterbreak\char`\#}{\char`\#}}%
            \def\PYGZpc{\discretionary{}{\Wrappedafterbreak\char`\%}{\char`\%}}%
            \def\PYGZdl{\discretionary{}{\Wrappedafterbreak\char`\$}{\char`\$}}%
            \def\PYGZhy{\discretionary{\char`\-}{\Wrappedafterbreak}{\char`\-}}%
            \def\PYGZsq{\discretionary{}{\Wrappedafterbreak\textquotesingle}{\textquotesingle}}%
            \def\PYGZdq{\discretionary{}{\Wrappedafterbreak\char`\"}{\char`\"}}%
            \def\PYGZti{\discretionary{\char`\~}{\Wrappedafterbreak}{\char`\~}}%
        }
        % Some characters . , ; ? ! / are not pygmentized.
        % This macro makes them "active" and they will insert potential linebreaks
        \newcommand*\Wrappedbreaksatpunct {%
            \lccode`\~`\.\lowercase{\def~}{\discretionary{\hbox{\char`\.}}{\Wrappedafterbreak}{\hbox{\char`\.}}}%
            \lccode`\~`\,\lowercase{\def~}{\discretionary{\hbox{\char`\,}}{\Wrappedafterbreak}{\hbox{\char`\,}}}%
            \lccode`\~`\;\lowercase{\def~}{\discretionary{\hbox{\char`\;}}{\Wrappedafterbreak}{\hbox{\char`\;}}}%
            \lccode`\~`\:\lowercase{\def~}{\discretionary{\hbox{\char`\:}}{\Wrappedafterbreak}{\hbox{\char`\:}}}%
            \lccode`\~`\?\lowercase{\def~}{\discretionary{\hbox{\char`\?}}{\Wrappedafterbreak}{\hbox{\char`\?}}}%
            \lccode`\~`\!\lowercase{\def~}{\discretionary{\hbox{\char`\!}}{\Wrappedafterbreak}{\hbox{\char`\!}}}%
            \lccode`\~`\/\lowercase{\def~}{\discretionary{\hbox{\char`\/}}{\Wrappedafterbreak}{\hbox{\char`\/}}}%
            \catcode`\.\active
            \catcode`\,\active
            \catcode`\;\active
            \catcode`\:\active
            \catcode`\?\active
            \catcode`\!\active
            \catcode`\/\active
            \lccode`\~`\~
        }
    \makeatother

    \let\OriginalVerbatim=\Verbatim
    \makeatletter
    \renewcommand{\Verbatim}[1][1]{%
        %\parskip\z@skip
        \sbox\Wrappedcontinuationbox {\Wrappedcontinuationsymbol}%
        \sbox\Wrappedvisiblespacebox {\FV@SetupFont\Wrappedvisiblespace}%
        \def\FancyVerbFormatLine ##1{\hsize\linewidth
            \vtop{\raggedright\hyphenpenalty\z@\exhyphenpenalty\z@
                \doublehyphendemerits\z@\finalhyphendemerits\z@
                \strut ##1\strut}%
        }%
        % If the linebreak is at a space, the latter will be displayed as visible
        % space at end of first line, and a continuation symbol starts next line.
        % Stretch/shrink are however usually zero for typewriter font.
        \def\FV@Space {%
            \nobreak\hskip\z@ plus\fontdimen3\font minus\fontdimen4\font
            \discretionary{\copy\Wrappedvisiblespacebox}{\Wrappedafterbreak}
            {\kern\fontdimen2\font}%
        }%

        % Allow breaks at special characters using \PYG... macros.
        \Wrappedbreaksatspecials
        % Breaks at punctuation characters . , ; ? ! and / need catcode=\active
        \OriginalVerbatim[#1,codes*=\Wrappedbreaksatpunct]%
    }
    \makeatother

    % Exact colors from NB
    \definecolor{incolor}{HTML}{303F9F}
    \definecolor{outcolor}{HTML}{D84315}
    \definecolor{cellborder}{HTML}{CFCFCF}
    \definecolor{cellbackground}{HTML}{F7F7F7}

    % prompt
    \makeatletter
    \newcommand{\boxspacing}{\kern\kvtcb@left@rule\kern\kvtcb@boxsep}
    \makeatother
    \newcommand{\prompt}[4]{
        {\ttfamily\llap{{\color{#2}[#3]:\hspace{3pt}#4}}\vspace{-\baselineskip}}
    }
    

    
    % Prevent overflowing lines due to hard-to-break entities
    \sloppy
    % Setup hyperref package
    \hypersetup{
      breaklinks=true,  % so long urls are correctly broken across lines
      colorlinks=true,
      urlcolor=urlcolor,
      linkcolor=linkcolor,
      citecolor=citecolor,
      }
    % Slightly bigger margins than the latex defaults
    
    \geometry{verbose,tmargin=1in,bmargin=1in,lmargin=1in,rmargin=1in}
    
    

\begin{document}
    
    \maketitle
    
    

    
    Probability for Data Science

UC Berkeley, Spring 2026

Ani Adhikari

CC BY-NC-SA 4.0

This content is protected and may not be shared, uploaded, or
distributed.

    \begin{tcolorbox}[breakable, size=fbox, boxrule=1pt, pad at break*=1mm,colback=cellbackground, colframe=cellborder]
\prompt{In}{incolor}{1}{\boxspacing}
\begin{Verbatim}[commandchars=\\\{\}]
\PY{k+kn}{import}\PY{+w}{ }\PY{n+nn}{warnings}
\PY{n}{warnings}\PY{o}{.}\PY{n}{filterwarnings}\PY{p}{(}\PY{l+s+s1}{\PYZsq{}}\PY{l+s+s1}{ignore}\PY{l+s+s1}{\PYZsq{}}\PY{p}{)}

\PY{k+kn}{from}\PY{+w}{ }\PY{n+nn}{prob140}\PY{+w}{ }\PY{k+kn}{import} \PY{o}{*}
\PY{k+kn}{from}\PY{+w}{ }\PY{n+nn}{datascience}\PY{+w}{ }\PY{k+kn}{import} \PY{o}{*}
\PY{k+kn}{import}\PY{+w}{ }\PY{n+nn}{numpy}\PY{+w}{ }\PY{k}{as}\PY{+w}{ }\PY{n+nn}{np}
\PY{k+kn}{from}\PY{+w}{ }\PY{n+nn}{scipy}\PY{+w}{ }\PY{k+kn}{import} \PY{n}{stats}

\PY{k+kn}{import}\PY{+w}{ }\PY{n+nn}{matplotlib}\PY{n+nn}{.}\PY{n+nn}{pyplot}\PY{+w}{ }\PY{k}{as}\PY{+w}{ }\PY{n+nn}{plt}
\PY{o}{\PYZpc{}}\PY{k}{matplotlib} inline
\PY{k+kn}{import}\PY{+w}{ }\PY{n+nn}{matplotlib}
\PY{n}{matplotlib}\PY{o}{.}\PY{n}{style}\PY{o}{.}\PY{n}{use}\PY{p}{(}\PY{l+s+s1}{\PYZsq{}}\PY{l+s+s1}{fivethirtyeight}\PY{l+s+s1}{\PYZsq{}}\PY{p}{)}
\end{Verbatim}
\end{tcolorbox}

    \hypertarget{section}{%
\section{}\label{section}}

    \hypertarget{homework-10-due-monday-april-13th-at-5pm}{%
\section{Homework 10 (Due Monday, April 13th at
5PM)}\label{homework-10-due-monday-april-13th-at-5pm}}

    \hypertarget{instructions}{%
\subsubsection{Instructions}\label{instructions}}

\hypertarget{questions-1-2-3-and-4-can-be-answered-based-on-tuesdays-lecture-and-past-content-alone.}{%
\paragraph{\texorpdfstring{\textbf{Questions 1, 2, 3, and 4 can be
answered based on Tuesday's lecture (and past content)
alone.}}{Questions 1, 2, 3, and 4 can be answered based on Tuesday's lecture (and past content) alone.}}\label{questions-1-2-3-and-4-can-be-answered-based-on-tuesdays-lecture-and-past-content-alone.}}

\hypertarget{we-strongly-encourage-you-to-attempt-questions-2-and-5-independently-relying-on-the-provided-textbook-sections-and-your-own-reasoning.-if-you-find-yourself-stuck-we-recommend-attending-office-hours-or-homework-party-or-posting-your-questions-on-ed.-if-youre-having-trouble-getting-started-on-problems-without-the-use-of-ai-we-encourage-you-to-follow-our-recommendation.}{%
\paragraph{\texorpdfstring{\textbf{We strongly encourage you to attempt
Questions 2 and 5 independently, relying on the provided textbook
sections and your own reasoning. If you find yourself stuck, we
recommend attending office hours or homework party or posting your
questions on Ed. If you're having trouble getting started on problems
without the use of AI, we encourage you to follow our
recommendation.}}{We strongly encourage you to attempt Questions 2 and 5 independently, relying on the provided textbook sections and your own reasoning. If you find yourself stuck, we recommend attending office hours or homework party or posting your questions on Ed. If you're having trouble getting started on problems without the use of AI, we encourage you to follow our recommendation.}}\label{we-strongly-encourage-you-to-attempt-questions-2-and-5-independently-relying-on-the-provided-textbook-sections-and-your-own-reasoning.-if-you-find-yourself-stuck-we-recommend-attending-office-hours-or-homework-party-or-posting-your-questions-on-ed.-if-youre-having-trouble-getting-started-on-problems-without-the-use-of-ai-we-encourage-you-to-follow-our-recommendation.}}

Your homeworks will generally have two components: a written portion and
a portion that also involves code. Written work should be completed on
paper, and coding questions should be done in the notebook. Start the
work for the written portions of each section on a new page. You are
welcome to \(\LaTeX\) your answers to the written portions, but staff
will not be able to assist you with \(\LaTeX\) related issues.

It is your responsibility to ensure that both components of the homework
are submitted completely and properly to Gradescope. \textbf{Make sure
to assign each page of your pdf to the correct question. Refer to the
bottom of the notebook for submission instructions.}

    \hypertarget{how-to-do-your-homework}{%
\subsubsection{How to Do Your Homework}\label{how-to-do-your-homework}}

The point of homework is for you to try your hand at using what you've
learned in class. The steps to follow:

\begin{itemize}
\tightlist
\item
  Go to lecture and sections, and also go over the relevant text
  sections before starting on the homework. This will remind you what
  was covered in class, and the text will typically contain examples not
  covered in lecture. The weekly Study Guide will list what you should
  read.
\item
  Work on some of the practice problems before starting on the homework.
\item
  Attempt the homework problems by yourself with the text, section work,
  and practice materials all at hand. Sometimes the week's lab will help
  as well. The two steps above will help this step go faster and be more
  fruitful.
\item
  At this point, seek help if you need it. Don't ask how to do the
  problem --- ask how to get started, or for a nudge to get you past
  where you are stuck. Always say what you have already tried. That
  helps us help you more effectively.
\item
  For a good measure of your understanding, keep track of the fraction
  of the homework you can do by yourself or with minimal help. It's a
  better measure than your homework score, and only you can measure it.
\end{itemize}

    \hypertarget{rules-for-homework}{%
\subsubsection{Rules for Homework}\label{rules-for-homework}}

\begin{itemize}
\tightlist
\item
  Every answer should contain a calculation or reasoning. For example, a
  calculation such as \((1/3)(0.8) + (2/3)(0.7)\) or
  \texttt{sum({[}(1/3)*0.8,\ (2/3)*0.7{]})}is fine without further
  explanation or simplification. If we want you to simplify, we'll ask
  you to. But just \({5 \choose 2}\) by itself is not fine; write ``we
  want any 2 out of the 5 frogs and they can appear in any order'' or
  whatever reasoning you used. Reasoning can be brief and abbreviated,
  e.g.~``product rule'' or ``not mutually exclusive.''
\item
  You may consult others (see ``How to Do Your Homework'' above) but you
  must write up your own answers using your own words, notation, and
  sequence of steps.
\item
  We'll be using Gradescope. You must submit the homework according to
  the instructions at the end of homework set.
\end{itemize}

\hypertarget{we-will-not-grade-assignments-which-do-not-have-pages-correctly-selected-for-each-question.}{%
\subsection{\texorpdfstring{\textbf{We will not grade assignments which
do not have pages correctly selected for each
question.}}{We will not grade assignments which do not have pages correctly selected for each question.}}\label{we-will-not-grade-assignments-which-do-not-have-pages-correctly-selected-for-each-question.}}

    \hypertarget{preliminary-line-plots}{%
\subsubsection{Preliminary: Line Plots}\label{preliminary-line-plots}}

In Data 8 you used \texttt{Table.plot} to draw line plots, but in this
class the line plots have typically been drawn using the plot function
of matplotlib. Here is the syntax; you are going to need it in the
exercises. It's easier than setting up tables first, especially when you
want to overlay multiple plots.

The \texttt{pyplot} module of matplotlib has been imported for you as
\texttt{plt}. This is its standard abbreviation.

Let \texttt{x} and \texttt{y} be two numerical arrays of the same
length. Then \texttt{plt.plot(x,\ y)} produces a line plot with values
of \texttt{x} on the horizontal axis and the corresponding values of
\texttt{y} on the vertical. The plot ``joins the dots''
\texttt{(x.item(0),\ y.item(0))}, \texttt{(x.item(1),\ y.item(1))},
\texttt{(x.item(2),\ y.item(2))}, and so on.

The optional argument \texttt{lw=n} can be used to set a line width of n
units. Please use \texttt{lw=2} in this homework.

The semicolon at the end of the last line of code in each cell stops
matplotlib from blurting out text that we don't need here.

Run these cells to see some examples. Note the overlaid plots: they are
straightforward to draw, and Python chooses a different color for each
plot.

    \begin{tcolorbox}[breakable, size=fbox, boxrule=1pt, pad at break*=1mm,colback=cellbackground, colframe=cellborder]
\prompt{In}{incolor}{3}{\boxspacing}
\begin{Verbatim}[commandchars=\\\{\}]
\PY{n}{x} \PY{o}{=} \PY{n}{np}\PY{o}{.}\PY{n}{arange}\PY{p}{(}\PY{l+m+mi}{0}\PY{p}{,} \PY{l+m+mf}{1.01}\PY{p}{,} \PY{l+m+mf}{0.01}\PY{p}{)}
\PY{n}{x\PYZus{}squared} \PY{o}{=} \PY{n}{x} \PY{o}{*}\PY{o}{*} \PY{l+m+mi}{2}
\PY{n}{x\PYZus{}cubed} \PY{o}{=} \PY{n}{x}\PY{o}{*}\PY{o}{*}\PY{l+m+mi}{3}
\PY{n}{plt}\PY{o}{.}\PY{n}{plot}\PY{p}{(}\PY{n}{x}\PY{p}{,} \PY{n}{x\PYZus{}squared}\PY{p}{,} \PY{n}{lw}\PY{o}{=}\PY{l+m+mi}{2}\PY{p}{)}
\PY{n}{plt}\PY{o}{.}\PY{n}{plot}\PY{p}{(}\PY{n}{x}\PY{p}{,} \PY{n}{x\PYZus{}cubed}\PY{p}{,} \PY{n}{lw}\PY{o}{=}\PY{l+m+mi}{2}\PY{p}{)}\PY{p}{;}
\end{Verbatim}
\end{tcolorbox}

    \begin{center}
    \end{center}
    { \hspace*{\fill} \\}
    
    You can add titles, labels, and legends.

    \begin{tcolorbox}[breakable, size=fbox, boxrule=1pt, pad at break*=1mm,colback=cellbackground, colframe=cellborder]
\prompt{In}{incolor}{4}{\boxspacing}
\begin{Verbatim}[commandchars=\\\{\}]
\PY{n}{plt}\PY{o}{.}\PY{n}{plot}\PY{p}{(}\PY{n}{x}\PY{p}{,} \PY{n}{x\PYZus{}squared}\PY{p}{,} \PY{n}{lw}\PY{o}{=}\PY{l+m+mi}{2}\PY{p}{,} \PY{n}{label} \PY{o}{=} \PY{l+s+s1}{\PYZsq{}}\PY{l+s+s1}{Square}\PY{l+s+s1}{\PYZsq{}}\PY{p}{)}
\PY{n}{plt}\PY{o}{.}\PY{n}{plot}\PY{p}{(}\PY{n}{x}\PY{p}{,} \PY{n}{x\PYZus{}cubed}\PY{p}{,} \PY{n}{lw}\PY{o}{=}\PY{l+m+mi}{2}\PY{p}{,} \PY{n}{label} \PY{o}{=} \PY{l+s+s1}{\PYZsq{}}\PY{l+s+s1}{Cube}\PY{l+s+s1}{\PYZsq{}}\PY{p}{)}
\PY{n}{plt}\PY{o}{.}\PY{n}{legend}\PY{p}{(}\PY{p}{)}
\PY{n}{plt}\PY{o}{.}\PY{n}{xlabel}\PY{p}{(}\PY{l+s+s1}{\PYZsq{}}\PY{l+s+s1}{x}\PY{l+s+s1}{\PYZsq{}}\PY{p}{)}
\PY{n}{plt}\PY{o}{.}\PY{n}{title}\PY{p}{(}\PY{l+s+s1}{\PYZsq{}}\PY{l+s+s1}{Powers of \PYZdl{}x\PYZdl{}}\PY{l+s+s1}{\PYZsq{}}\PY{p}{)}\PY{p}{;}
\end{Verbatim}
\end{tcolorbox}

    \begin{center}
    \end{center}
    { \hspace*{\fill} \\}
    
    As another example, the code below was used in Homework 8 Exercise 4d to
plot gamma densities.

    \begin{tcolorbox}[breakable, size=fbox, boxrule=1pt, pad at break*=1mm,colback=cellbackground, colframe=cellborder]
\prompt{In}{incolor}{5}{\boxspacing}
\begin{Verbatim}[commandchars=\\\{\}]
\PY{n}{t} \PY{o}{=} \PY{n}{np}\PY{o}{.}\PY{n}{arange}\PY{p}{(}\PY{l+m+mi}{0}\PY{p}{,} \PY{l+m+mi}{25}\PY{p}{,} \PY{l+m+mf}{0.01}\PY{p}{)}
\PY{n}{r\PYZus{}1} \PY{o}{=} \PY{n}{stats}\PY{o}{.}\PY{n}{gamma}\PY{o}{.}\PY{n}{pdf}\PY{p}{(}\PY{n}{t}\PY{p}{,} \PY{l+m+mi}{1}\PY{p}{,} \PY{n}{scale}\PY{o}{=}\PY{l+m+mi}{1}\PY{o}{/}\PY{l+m+mf}{0.25}\PY{p}{)}
\PY{n}{r\PYZus{}1\PYZus{}5} \PY{o}{=} \PY{n}{stats}\PY{o}{.}\PY{n}{gamma}\PY{o}{.}\PY{n}{pdf}\PY{p}{(}\PY{n}{t}\PY{p}{,} \PY{l+m+mf}{1.5}\PY{p}{,} \PY{n}{scale}\PY{o}{=}\PY{l+m+mi}{1}\PY{o}{/}\PY{l+m+mf}{0.25}\PY{p}{)}
\PY{n}{r\PYZus{}2} \PY{o}{=} \PY{n}{stats}\PY{o}{.}\PY{n}{gamma}\PY{o}{.}\PY{n}{pdf}\PY{p}{(}\PY{n}{t}\PY{p}{,} \PY{l+m+mi}{2}\PY{p}{,} \PY{n}{scale}\PY{o}{=}\PY{l+m+mi}{1}\PY{o}{/}\PY{l+m+mf}{0.25}\PY{p}{)}
\PY{n}{r\PYZus{}2\PYZus{}5} \PY{o}{=} \PY{n}{stats}\PY{o}{.}\PY{n}{gamma}\PY{o}{.}\PY{n}{pdf}\PY{p}{(}\PY{n}{t}\PY{p}{,} \PY{l+m+mf}{2.5}\PY{p}{,} \PY{n}{scale}\PY{o}{=}\PY{l+m+mi}{1}\PY{o}{/}\PY{l+m+mf}{0.25}\PY{p}{)}
\PY{n}{r\PYZus{}3} \PY{o}{=} \PY{n}{stats}\PY{o}{.}\PY{n}{gamma}\PY{o}{.}\PY{n}{pdf}\PY{p}{(}\PY{n}{t}\PY{p}{,} \PY{l+m+mi}{3}\PY{p}{,} \PY{n}{scale}\PY{o}{=}\PY{l+m+mi}{1}\PY{o}{/}\PY{l+m+mf}{0.25}\PY{p}{)}
\PY{n}{plt}\PY{o}{.}\PY{n}{plot}\PY{p}{(}\PY{n}{t}\PY{p}{,} \PY{n}{r\PYZus{}1}\PY{p}{,} \PY{n}{lw}\PY{o}{=}\PY{l+m+mi}{2}\PY{p}{,} \PY{n}{label}\PY{o}{=}\PY{l+s+s1}{\PYZsq{}}\PY{l+s+s1}{gamma (1, 0.25)}\PY{l+s+s1}{\PYZsq{}}\PY{p}{)}
\PY{n}{plt}\PY{o}{.}\PY{n}{plot}\PY{p}{(}\PY{n}{t}\PY{p}{,} \PY{n}{r\PYZus{}1\PYZus{}5}\PY{p}{,} \PY{n}{lw}\PY{o}{=}\PY{l+m+mi}{2}\PY{p}{,} \PY{n}{label}\PY{o}{=}\PY{l+s+s1}{\PYZsq{}}\PY{l+s+s1}{gamma (1.5, 0.25)}\PY{l+s+s1}{\PYZsq{}}\PY{p}{)}
\PY{n}{plt}\PY{o}{.}\PY{n}{plot}\PY{p}{(}\PY{n}{t}\PY{p}{,} \PY{n}{r\PYZus{}2}\PY{p}{,} \PY{n}{lw}\PY{o}{=}\PY{l+m+mi}{2}\PY{p}{,} \PY{n}{label}\PY{o}{=}\PY{l+s+s1}{\PYZsq{}}\PY{l+s+s1}{gamma (2, 0.25)}\PY{l+s+s1}{\PYZsq{}}\PY{p}{)}
\PY{n}{plt}\PY{o}{.}\PY{n}{plot}\PY{p}{(}\PY{n}{t}\PY{p}{,} \PY{n}{r\PYZus{}2\PYZus{}5}\PY{p}{,} \PY{n}{lw}\PY{o}{=}\PY{l+m+mi}{2}\PY{p}{,} \PY{n}{label}\PY{o}{=}\PY{l+s+s1}{\PYZsq{}}\PY{l+s+s1}{gamma (2.5, 0.25)}\PY{l+s+s1}{\PYZsq{}}\PY{p}{)}
\PY{n}{plt}\PY{o}{.}\PY{n}{plot}\PY{p}{(}\PY{n}{t}\PY{p}{,} \PY{n}{r\PYZus{}3}\PY{p}{,} \PY{n}{lw}\PY{o}{=}\PY{l+m+mi}{2}\PY{p}{,} \PY{n}{label}\PY{o}{=}\PY{l+s+s1}{\PYZsq{}}\PY{l+s+s1}{gamma (3, 0.25)}\PY{l+s+s1}{\PYZsq{}}\PY{p}{)}
\PY{n}{plt}\PY{o}{.}\PY{n}{xlabel}\PY{p}{(}\PY{l+s+s1}{\PYZsq{}}\PY{l+s+s1}{\PYZdl{}t\PYZdl{}}\PY{l+s+s1}{\PYZsq{}}\PY{p}{)}
\PY{n}{plt}\PY{o}{.}\PY{n}{legend}\PY{p}{(}\PY{p}{)}
\PY{n}{plt}\PY{o}{.}\PY{n}{title}\PY{p}{(}\PY{l+s+s1}{\PYZsq{}}\PY{l+s+s1}{Gamma Densities: Same Rate and Different Shape Parameters}\PY{l+s+s1}{\PYZsq{}}\PY{p}{)}\PY{p}{;}
\end{Verbatim}
\end{tcolorbox}

    \begin{center}
    \end{center}
    { \hspace*{\fill} \\}
    
    \newpage

    \hypertarget{poisson-mgf}{%
\subsection{1. Poisson MGF}\label{poisson-mgf}}

Let \(X\) have Poisson(\(\mu\)) distribution, and let \(Y\) independent
of \(X\) have Poisson \((\lambda)\) distribution.

\textbf{a)} Derive the moment generating function of \(X\). The answer
is in the textbook. Your task here is to show why it's true.

\begin{align*}
	M_{X}(t) &= E[e^{tX}] \\
	&= \sum_{k=0}^{\infty } e^{tk} e^{-\mu } \frac{\mu ^{k}}{k!} \\
	&= e^{-\mu } \sum_{k=0}^{\infty } e^{tk} \frac{\mu ^{k}}{k!} \\
	&= e^{-\mu } \left( 1 + e^{t} \frac{\mu }{1} + e^{2t} \frac{\mu ^{2}}{2!} + \dots\right)  \\
	&= e^{-\mu } \left( 1 + \frac{\left( e^{t}\mu  \right)^{1} }{1!} + \frac{\left( e^{t}\mu  \right) ^{2}}{2!} + \dots \right)  \\
	&= e^{-\mu } e^{e^{t}\mu } \\
	&= e^{e^{t}\mu  - \mu } \\
	&= e^{\mu ( e^{t} - 1)} \\
\end{align*}



\textbf{b)} Use the result of (a) to show that the distribution of
\(X+Y\) is Poisson.


\begin{align*}
	M_{X+Y}(t) &= E[e^{tX}e^{tY}] \\
	&= e^{\mu (e^{t}-1)} \cdot e^{\lambda (e^{t}-1)} \\
	&= e^{\mu (e^{t}-1) + \lambda(e^{t-1)})} \\
	&= e^{(e^{t}-1)(\mu +\lambda)} \\
\end{align*}

This is the MGF for a Poisson distribution with parameter $\mu + \lambda$ 


    \newpage

    \hypertarget{gamma-tail-bound}{%
\subsection{2. Gamma Tail Bound}\label{gamma-tail-bound}}

Before you do this exercise, carefully study a
\href{https://data140.org/textbook/content/chapter-19/chernoff-bound/\#application-to-the-normal-distribution}{relevant
example} in the textbook. You will have to follow similar steps.

You will need the
\href{https://data140.org/textbook/content/chapter-19/moment-generating-functions/\#mgf-of-a-gamma-r-lambda-random-variable}{mgf
of the gamma distribution}. Also remember that you found the gamma mean
and variance in Homework 8.

If you want more information, read through
\href{https://data140.org/textbook/content/chapter-12/bounds/}{Ch. 12.3}
(Tail Bounds),
\href{https://data140.org/textbook/content/chapter-19/moment-generating-functions/}{Ch.
19.2} (MGFs), and
\href{https://data140.org/textbook/content/chapter-19/chernoff-bound/}{Ch.
19.4} (Chernoff Bounds) in its entirety. If you have any questions, come
to office hours or post them on Ed.

Let \(X\) have the gamma \((r, \lambda)\) distribution.

\textbf{a)} Show that \(P(X \ge 2E(X)) \le \left(\frac{2}{e}\right)^r\).


\begin{align*}
	P(X \ge 2E[X]) &= P(e^{tX} \ge e^{2tE[X]})  \\
		       &\le \min_{t>0}\frac{E[e^{tX}]}{e^{2tE[X]}} \\
		       &= \min_{t>0} \frac{M_{X}(t)}{e^{2tE[X]}} \\
		       &= \min_{t>0}\left( \frac{\lambda}{\lambda-t} \right) ^{r} / e^{2 \frac{rt}{\lambda}}  \\
\end{align*}

The minimum value of this function is where the log of that function is minimized.

$$
	f(t) = \log (\left( \frac{\lambda}{\lambda - rt} \right) ^{r} / e^{2 \frac{rt}{\lambda}}) = r\log\left( \frac{\lambda}{\lambda-t} \right) - 2 \frac{rt}{\lambda} = r\log(\lambda) - r\log(\lambda-t) - \frac{2rt}{\lambda}\\
$$

\begin{align*}
	\frac{df}{dt} &= \frac{r}{\lambda - \hat t} - \frac{2r}{\lambda} = 0\\
	\lambda - \hat t &= \frac{\lambda}{2} \\
	\hat t &= \lambda - \frac{\lambda}{2}  \\
	\hat t &= \frac{\lambda}{2} \\
\end{align*}

Plugging back into the original function:

\begin{align*}
	\min_{t>0} \left( \frac{\lambda}{\lambda-t}^{r} \right) / e^{2 \frac{rt}{\lambda}} &= \left( \frac{\lambda}{\lambda - \frac{\lambda}{2}} \right)^{r} / e^{2 r \frac{\lambda}{2} / \lambda} \\
	&= \left( \frac{\lambda}{\frac{\lambda}{2}} \right) ^{r} / e^{r} \\
	&= \left( \frac{2^{r}}{e^{r}} \right)  \\
	&= \left( \frac{2}{e} \right) ^{r} \\
\end{align*}


\textbf{b)} Find Markov's and Chebyshev's bounds on \(P(X \ge 2E(X))\).


Markov:
\begin{align*}
	P(X \ge 2E[X]) &\le \frac{E[X]}{2E[X]}\\
	&= \frac{1}{2} \\
\end{align*}

Chebyshev's

\begin{align*}
	P(X \ge 2E[X]) &= P(X - E[X] \ge E[X]) \\
		       &\le P(\vert X - E[X] \vert \ge E[X]) \\
		       &\le \frac{Var(X)}{E[X]^{2}} \\ 
		       &\le \frac{r}{\lambda^{2}} / \frac{r^{2}}{\lambda^{2}}\\
		       &= \frac{1}{r} \\
\end{align*}


\textbf{c) {[}CODE{]}} Fix \(\lambda = 1\). Display overlaid plots of
the following four graphs as functions of \(r\), for \(r\) in the
interval \((0.5, 15)\) :

\begin{itemize}
\tightlist
\item
  The exact tail probability \(P(X \ge 2E(X))\)
\item
  The bound in Part \textbf{a}: \(\left(\frac{2}{e}\right)^r\)
\item
  Chebyshev's bound on \(P(X \ge 2E(X))\)
\item
  Markov's bound on \(P(X \ge 2E(X))\)
\end{itemize}

The code uses \texttt{plt.plot} which you have used before. The
expression \texttt{stats.gamma.cdf(x,\ r,\ scale=1)} evaluates to the
cdf of the gamma \((r, 1)\) distribution at the point \(x\).

    \begin{tcolorbox}[breakable, size=fbox, boxrule=1pt, pad at break*=1mm,colback=cellbackground, colframe=cellborder]
\prompt{In}{incolor}{ }{\boxspacing}
\begin{Verbatim}[commandchars=\\\{\}]
\PY{c+c1}{\PYZsh{} Answer to c}
\PY{n}{r} \PY{o}{=} \PY{n}{np}\PY{o}{.}\PY{n}{arange}\PY{p}{(}\PY{l+m+mf}{0.05}\PY{p}{,} \PY{l+m+mi}{15}\PY{p}{,} \PY{l+m+mf}{0.1}\PY{p}{)} 

\PY{n}{markov\PYZus{}bound} \PY{o}{=} \PY{o}{.}\PY{o}{.}\PY{o}{.}

\PY{n}{chebyshev\PYZus{}bound} \PY{o}{=} \PY{o}{.}\PY{o}{.}\PY{o}{.}

\PY{n}{part\PYZus{}a\PYZus{}bound} \PY{o}{=} \PY{o}{.}\PY{o}{.}\PY{o}{.}

\PY{c+c1}{\PYZsh{} Use as many lines as you need for the exact values}
\PY{n}{exact} \PY{o}{=} \PY{o}{.}\PY{o}{.}\PY{o}{.}
\PY{o}{.}\PY{o}{.}\PY{o}{.}

\PY{n}{plt}\PY{o}{.}\PY{n}{plot}\PY{p}{(}\PY{n}{r}\PY{p}{,} \PY{n}{exact}\PY{p}{,} \PY{n}{lw}\PY{o}{=}\PY{l+m+mi}{2}\PY{p}{,} \PY{n}{label}\PY{o}{=}\PY{l+s+s1}{\PYZsq{}}\PY{l+s+s1}{Exact Chance}\PY{l+s+s1}{\PYZsq{}}\PY{p}{)}
\PY{n}{plt}\PY{o}{.}\PY{n}{plot}\PY{p}{(}\PY{n}{r}\PY{p}{,} \PY{n}{part\PYZus{}a\PYZus{}bound}\PY{p}{,} \PY{n}{lw}\PY{o}{=}\PY{l+m+mi}{2}\PY{p}{,} \PY{n}{label}\PY{o}{=}\PY{l+s+s1}{\PYZsq{}}\PY{l+s+s1}{Part (a) Bound}\PY{l+s+s1}{\PYZsq{}}\PY{p}{)}
\PY{n}{plt}\PY{o}{.}\PY{n}{plot}\PY{p}{(}\PY{n}{r}\PY{p}{,} \PY{n}{chebyshev\PYZus{}bound}\PY{p}{,} \PY{n}{lw}\PY{o}{=}\PY{l+m+mi}{2}\PY{p}{,} \PY{n}{label}\PY{o}{=}\PY{l+s+s1}{\PYZsq{}}\PY{l+s+s1}{Chebyshev Bound}\PY{l+s+s1}{\PYZsq{}}\PY{p}{)}
\PY{n}{plt}\PY{o}{.}\PY{n}{plot}\PY{p}{(}\PY{n}{r}\PY{p}{,} \PY{n}{markov\PYZus{}bound}\PY{p}{,} \PY{n}{lw}\PY{o}{=}\PY{l+m+mi}{2}\PY{p}{,} \PY{n}{label}\PY{o}{=}\PY{l+s+s1}{\PYZsq{}}\PY{l+s+s1}{Markov Bound}\PY{l+s+s1}{\PYZsq{}}\PY{p}{)}
\PY{n}{plt}\PY{o}{.}\PY{n}{legend}\PY{p}{(}\PY{p}{)}
\PY{n}{plt}\PY{o}{.}\PY{n}{xlabel}\PY{p}{(}\PY{l+s+s1}{\PYZsq{}}\PY{l+s+s1}{\PYZdl{}r\PYZdl{}}\PY{l+s+s1}{\PYZsq{}}\PY{p}{)}
\PY{n}{plt}\PY{o}{.}\PY{n}{ylim}\PY{p}{(}\PY{l+m+mi}{0}\PY{p}{,} \PY{l+m+mi}{1}\PY{p}{)}
\PY{n}{plt}\PY{o}{.}\PY{n}{xlim}\PY{p}{(}\PY{l+m+mi}{0}\PY{p}{,} \PY{l+m+mi}{15}\PY{p}{)}
\PY{n}{plt}\PY{o}{.}\PY{n}{title}\PY{p}{(}\PY{l+s+s1}{\PYZsq{}}\PY{l+s+s1}{\PYZdl{}P(X }\PY{l+s+s1}{\PYZbs{}}\PY{l+s+s1}{geq 2E(X))\PYZdl{} for \PYZdl{}X\PYZdl{} gamma \PYZdl{}(r, 1)\PYZdl{}}\PY{l+s+s1}{\PYZsq{}}\PY{p}{)}\PY{p}{;}
\end{Verbatim}
\end{tcolorbox}

    \newpage

    \hypertarget{mle-of-the-exponential-rate}{%
\subsection{3. MLE of the Exponential
Rate}\label{mle-of-the-exponential-rate}}

For \(n > 1\), let \(X_1, X_2, \ldots , X_n\) be i.i.d. exponential
\((\lambda )\) variables.

\textbf{a)} Let \(\hat{\lambda}_n\) be the maximum likelihood estimate
(MLE) of the parameter \(\lambda\). Find \(\hat{\lambda}_n\) in terms of
the sample mean \(\bar{X}_n = \frac{1}{n} \sum_{i=1}^n X_i\). The
subscript \(n\) in \(\bar{X}_n\) is there to remind us that we have the
average of \(n\) values. It doesn't refer to the \(n\)th sampled value
\(X_n\).


\begin{align*}
	lik(\lambda) &= \prod_{i=1}^{n}\lambda e^{-\lambda X_{i}} \\
	&= \lambda^{n} e^{-n\lambda\sum_{i=1}^{n}X_{i}}\\
\end{align*}

$$ \min_{\lambda} lik(\lambda) = \min_{\lambda} L(\lambda)$$ 

\begin{align*}
	L(\lambda) &=\log( \lambda^{n}e^{-\lambda \sum_{i=1}^{n} X_{i}}) \\
&= n\log(\lambda) - \lambda \sum_{i=1}^{n}X_{i} \\
\end{align*}


\begin{align*}
\frac{dL}{d\lambda} &= \frac{n}{\lambda} -  \sum_{i=1}^{n}X_{i} = 0\\ 
\frac{n}{\lambda} &= \sum_{i=1}^{n}X_{i} \\
\hat \lambda &=  \frac{n}{\sum_{i=1}^{n}X_{i}}  \\
&= \frac{1}{\bar X_{n}} \\
\end{align*}



\textbf{b)} Use facts about sums and linear transformations to find the
distribution of \(\bar{X}_n\) with little or no calculation. Recognize
it as one of the famous ones and provide its name and parameters. Use it
to find \(E(\hat{\lambda}_n)\).

The sum of exponential random variables is a gamma distribution. Then we apply a linear transformation of $\frac{1}{n}$ to that gamma distribution. Thus, $\sum_{i=1}^{n}X$ is a Gamma($n, \lambda$), and $\bar X$ is a Gamma($n, n\lambda$) distribution. Thinking about it as adding i.i.d. exponential($n\lambda$) r.v.s also works.  


\begin{align*}
	E[\hat \lambda_{n}] &= E[\frac{1}{\bar X_{n}}]\\
	&= \int_{0}^{\infty } \frac{1}{x} f_{\bar X_{n}}(x) dx \\
	&= \int_{0}^{\infty } x^{-1} \frac{(n\lambda)^{n}}{\Gamma(n)} x^{n-1}e^{-n\lambda x}dx \\
	&= \int_{0}^{\infty } \frac{(n\lambda)^{n}}{\Gamma(n)} x^{(n-1) - 1}e^{-n\lambda x}dx \\
	&= \frac{(n\lambda)^{n}}{\Gamma(n)} \int_{0}^{\infty } x^{(n-1) -1}e^{-n\lambda x}dx \\
	&= \frac{(n\lambda)^{n}}{\Gamma(n)} \frac{\Gamma(n-1)}{(n\lambda)^{n-1}} \\
	&= n\lambda \cdot \frac{\Gamma(n-1)}{\Gamma(n)} \\
	&= \frac{n\lambda}{n-1} \\
	&= \lambda \left(\frac{n}{n-1}\right) \\
\end{align*}




\textbf{c)} Is \(\hat{\lambda}_n\) an unbiased estimate of \(\lambda\)?
If it is biased, does it overestimate on average, or does it
underestimate? Is it asymptotically unbiased? That is, does
\(E(\hat{\lambda}_n)\) converge to \(\lambda\) as \(n \to \infty\)?

No, $\hat \lambda$ is not an unbiased estimator of $\lambda$. It will overestimate on average, because of the $\frac{n}{n-1}$ term. However, as $n\to \infty $, $\frac{n}{n-1} \to 1$, which means that it converges to $\lambda$. 
    \newpage

    \hypertarget{mle-and-map-estimates}{%
\subsection{4. MLE and MAP Estimates}\label{mle-and-map-estimates}}

The coin is tossed 10 times and the resulting sequence is HTTHHHTHTH. In
the parts below, we refer to this information as ``the data''.

\textbf{a)} Under the assumption that the coin lands heads with a fixed
unknown probability \(p\), find the MLE of \(p\) based on the data.


\begin{align*}
	lik(p) &= p^{6}(1-p)^{4} \\ 
	L(p) &= 6\log(p) + 4\log(1-p)\\
	\frac{dL}{dp} &= \frac{6}{p} + \frac{4}{1-p} \\
	\frac{6}{\hat p} - \frac{4}{1 - \hat p} &= 0 \\
	\frac{6(1-\hat p) - 4p}{1-\hat p } &= 0 \\
	6 - 6\hat p - 4 \hat p  &= 0 \\
	10\hat p &= 6 \\
\hat p = \frac{6}{10} 
\end{align*}


\textbf{b)} Suppose now that the coin lands heads with a random
probability \(X\). Let the prior density of \(X\) be uniform on the unit
interval. Find the MAP estimate of the probability of heads, given the
data.


We showed during lecture that the MAP estimate given a prior beta(r, s) density, will be beta(r+k, s+n-k).

A uniform(0,1) distribution is the same as the beta(1,1) distribution. Thus, the postier distribution will be beta(1+6 = 7, 1+4 =5). Therefore the mode of this distribution is:

$$\frac{r+k-1}{r+s+n-2} = \frac{6}{10}$$

\textbf{c)} Show that if \(r > 1\) and \(s > 1\) then the mode of the
beta \((r, s)\) distribution is \((r-1)/(r+s-2)\). Remember to ignore
multiplicative constants and take the log before maximizing.


The beta pdf is:

$$\frac{\Gamma(r+s)}{\Gamma(r)\Gamma(s)}x^{r-1}(1-x)^{s-1} = C(r,s)x^{r-1}(1-x)^{s-1}$$

The mode is where this density is maximised. Since r and s are constants, we only care about the x value, and we do not care about the constant as it does not depend on x. 
Furthermore, maximizing the log of this function also maximises the original. 

Thus
$$
	\max_{x, r>1, s>1} g(x)  = (r-1)\log(x) + (s-1)\log(1-x) 
$$

\begin{align*}
	\frac{dg}{dx} = \frac{r-1}{x} - \frac{s-1}{1-x} &= 0 \\
	(r-1)(1-x) - (s-1)(x) &= 0 \\
	(r - rx -1 +x) - (sx - x) &= 0 \\
	r - rx - sx - 1 +2x &= 0 \\
	r - 1 -x(r + s - 2) &= 0 \\
	\boxed{\hat x = \frac{1-r}{r+s-2}}
\end{align*}

\textbf{d)} Suppose instead that the prior density of \(X\) is
\(f(x) = 4x^3\) if \(0 < x < 1\) and \(0\) otherwise. Find the MAP
estimate of the probability of heads, given the data.


\begin{align*}
	f_{X \mid S_{10}=  6} &\propto f_{X}(p)dp \binom{10}{6}p^{6}(1-p)^{4} \\
	&= 4p^{3} \binom{10}{6} p^{6}(1-p)^{4} \\
	&\propto p^{3} p^{6}(1-p)^{4} = p^{9}(1-p)^{4}
\end{align*}

\begin{align*}
	\log(f(p)) &= 9\log(p) + 4\log(1-p) \\
	\frac{9}{\hat p} - \frac{4}{1-\hat p } &= 0 \\
	9(1-\hat p) - 4(\hat p) &= 0\\
	9 - 13\hat p   &= 0 \\
	\hat p &= \frac{9}{13} \\
\end{align*}

    \newpage

    \hypertarget{heads-in-tosses-of-a-random-coin}{%
\subsection{5. Heads in Tosses of a Random
Coin}\label{heads-in-tosses-of-a-random-coin}}

Before you do this exercise, carefully read through
\href{https://data140.org/textbook/content/chapter-20/independence-revisited/}{Ch.
20.2} (Independence, Revisisted) and
\href{https://data140.org/textbook/content/chapter-21/updating-and-prediction/}{Ch.
21.1} (Updating and Prediction). For more assistance, look through Ch.
21 Ex.2 and Ch. 21 Ex.3, which is done in Friday's section. If you have
any questions, come to office hours or post them on Ed.

Let \(X\) be a random proportion with a prior distribution that is beta
\((2, 3)\). Given \(X=p\), let \(I_1, I_2, I_3, \ldots\) be i.i.d.
Bernoulli \((p)\).

\textbf{a)} Plot the prior density of \(X\).

    \begin{tcolorbox}[breakable, size=fbox, boxrule=1pt, pad at break*=1mm,colback=cellbackground, colframe=cellborder]
\prompt{In}{incolor}{ }{\boxspacing}
\begin{Verbatim}[commandchars=\\\{\}]
\PY{c+c1}{\PYZsh{} Answer to a}

\PY{n}{x} \PY{o}{=} \PY{n}{np}\PY{o}{.}\PY{n}{arange}\PY{p}{(}\PY{l+m+mi}{0}\PY{p}{,} \PY{l+m+mf}{1.01}\PY{p}{,} \PY{l+m+mf}{0.01}\PY{p}{)}
\PY{o}{.}\PY{o}{.}\PY{o}{.}
\PY{n}{plt}\PY{o}{.}\PY{n}{title}\PY{p}{(}\PY{l+s+s1}{\PYZsq{}}\PY{l+s+s1}{Beta \PYZdl{}(2, 3)\PYZdl{} density}\PY{l+s+s1}{\PYZsq{}}\PY{p}{)}\PY{p}{;}
\end{Verbatim}
\end{tcolorbox}

    \textbf{b)} Let \(H_7 = \sum_{k=1}^7 I_k\). For each
\(h = 0, 1, \ldots, 7\), plot the posterior density of \(X\) given
\(H_7 = h\). All \(8\) plots should be on the same graph.

Use as many lines of code as you need. You don't have to include labels
and a legend, but the title should say which densities you are plotting.

    \begin{tcolorbox}[breakable, size=fbox, boxrule=1pt, pad at break*=1mm,colback=cellbackground, colframe=cellborder]
\prompt{In}{incolor}{ }{\boxspacing}
\begin{Verbatim}[commandchars=\\\{\}]
\PY{c+c1}{\PYZsh{} Answer to b}

\PY{o}{.}\PY{o}{.}\PY{o}{.}

\PY{n}{plt}\PY{o}{.}\PY{n}{title}\PY{p}{(}\PY{l+s+s1}{\PYZsq{}}\PY{l+s+s1}{... (...) Densities for \PYZdl{}0 }\PY{l+s+s1}{\PYZbs{}}\PY{l+s+s1}{leq h }\PY{l+s+s1}{\PYZbs{}}\PY{l+s+s1}{leq 7}\PY{l+s+s1}{\PYZsq{}}\PY{p}{)}\PY{p}{;}
\end{Verbatim}
\end{tcolorbox}

    \textbf{c)} What is the MAP estimate of the random probability of heads
given \(H_7 = 2\)? Calculate the estimate using the appropriate formula
and confirm that your answer agrees with the estimate visible in the
appropriate graph above.

\textbf{d)} Find \(P(I_8 = 1 \mid H_7 = 2)\). Your answer should be a
decimal value.

\textbf{e)} Find \(P(I_8 = 1, I_9 = 1, I_{10} = 1 \mid H_7 = 2)\). Your
answer should be a decimal value. Is it equal to the cube of the answer
to Part (d)? If not, which is bigger?

    \newpage

    \hypertarget{submission-instructions}{%
\subsection{Submission Instructions}\label{submission-instructions}}

Many assignments throughout the course will have a written portion and a
code portion. Please follow the directions below to properly submit both
portions.

\hypertarget{written-portion}{%
\subsubsection{Written Portion}\label{written-portion}}

\begin{itemize}
\tightlist
\item
  Scan all the pages into a PDF. You can use any scanner or a phone
  using applications such as CamScanner. Please \textbf{DO NOT} simply
  take pictures using your phone.
\item
  Please start a new page for each question. If you have already written
  multiple questions on the same page, you can crop the image in
  CamScanner or fold your page over (the old-fashioned way). This helps
  expedite grading.
\item
  It is your responsibility to check that all the work on all the
  scanned pages is legible.
\item
  If you used \(\LaTeX\) to do the written portions, you do not need to
  do any scanning; you can just download the whole notebook as a PDF via
  LaTeX.
\end{itemize}

\hypertarget{code-portion}{%
\subsubsection{Code Portion}\label{code-portion}}

\begin{itemize}
\tightlist
\item
  Save your notebook using
  \texttt{File\ \textgreater{}\ Save\ and\ Checkpoint}.
\item
  Generate a PDF file using
  \texttt{File\ \textgreater{}\ Download\ As\ \textgreater{}\ PDF\ via\ LaTeX}.
  This might take a few seconds and will automatically download a PDF
  version of this notebook.

  \begin{itemize}
  \tightlist
  \item
    If you have issues, please post a follow-up on the general Homework
    10 Ed thread.
  \end{itemize}
\end{itemize}

\hypertarget{submitting}{%
\subsubsection{Submitting}\label{submitting}}

\begin{itemize}
\tightlist
\item
  Combine the PDFs from the written and code portions into one PDF.
  \href{https://smallpdf.com/merge-pdf}{Here} is a useful tool for doing
  so.
\item
  Submit the assignment to Homework 10 on Gradescope.
\item
  \textbf{Make sure to assign each page of your pdf to the correct
  question.}
\item
  \textbf{It is your responsibility to verify that all of your work
  shows up in your final PDF submission.}
\end{itemize}

If you are having difficulties scanning, uploading, or submitting your
work, please read the
\href{https://edstem.org/us/courses/94054/discussion/7533569}{Ed Thread}
on this topic and post a follow-up on the general Homework 10 Ed thread.


    % Add a bibliography block to the postdoc
    
    
    
\end{document}
